\chapter{Vertex Colourings}
%%%%%%%%%%%%%%%%%%%%%%%%%%%%%%%%%%%%%%%%%%%%%%%%%%%%%%%%%%%%%%%%%%%%%%%%%%%%%%%
\section{Vertex Colourings}

% Generated by scripts/blueprint_restate.py from TCSlib.GraphTheory.Core.VertexColourings.
% Each body is the STATEMENT only, taken from the declaration's Lean docstring:
% the one-line summary plus the verbatim book statement.  Proof, proof plan,
% significance commentary and formalisation notes are excluded by construction.

\begin{definition}[$G$ is critical: $\chi(H) < \chi(G)$ for every proper subgraph $H \subset G$]
\lean{SimpleGraph.IsCritical}
\label{SimpleGraph.IsCritical}
\leanok
\texttt{G} is \textbf{critical}: \texttt{$\chi$(H) < $\chi$(G)} for every proper subgraph \texttt{H $\subset$ G}.

\emph{It is helpful, when dealing with
colourings, to study the properties of a special class of graphs called critical
graphs. We say that a graph \texttt{G} is critical if \texttt{$\chi$(H) < $\chi$(G)} for
every proper
subgraph \texttt{H} of \texttt{G}. Such graphs were first investigated by Dirac (1952).}
\end{definition}

\begin{definition}[$G$ is $k$-critical: $k$-chromatic and critical]
\lean{SimpleGraph.IsKCritical}
\label{SimpleGraph.IsKCritical}
\leanok
\uses{SimpleGraph.IsCritical}
\texttt{G} is \texttt{k}-\textbf{critical}: \texttt{k}-chromatic and critical.

\emph{A \texttt{k}-critical graph is one that is
\texttt{k}-chromatic and critical; every \texttt{k}-chromatic graph has a \texttt{k}-critical
subgraph.  A 4-critical graph, due to Gr\"otzsch (1958), is shown in figure 8.2.}
\end{definition}

\begin{definition}[$G \cdot uv$: contract the edge $uv$ by identifying $v$ into $u$]
\lean{SimpleGraph.contractEdge}
\label{SimpleGraph.contractEdge}
\leanok
\texttt{G $\cdot$ uv}: contract the edge \texttt{uv} by identifying \texttt{v} into \texttt{u}. Carrier drops to \texttt{\{x // x $\ne$ v\}}.

\emph{An edge \texttt{e} of
\texttt{G} is said to be contracted if it is deleted and its ends are identified; the
resulting graph is denoted by \texttt{G $\cdot$ e}.}

Theorem 8.3(ii) restates it locally: \emph{\texttt{G$_{2}$ $\cdot$ uv} denotes the graph
obtained from
\texttt{G$_{2}$} by identifying \texttt{u} and \texttt{v}.}
\end{definition}

\begin{definition}[$\pi_k(G)$, the number of distinct proper $k$-colourings]
\lean{SimpleGraph.numColorings}
\label{SimpleGraph.numColorings}
\leanok
\texttt{$\pi$\_k(G)}, the number of distinct proper \texttt{k}-colourings.

\emph{We shall denote the number of distinct
\texttt{k}-colourings of \texttt{G} by \texttt{$\pi$\_k(G)}; thus \texttt{$\pi$\_k(G) > 0} if and only if \texttt{G} is
\texttt{k}-colourable.  Two colourings are to be regarded as distinct if some vertex is
assigned different colours in the two colourings; in other words, if
\texttt{(V$_{1}$, V$_{2}$, \dots{}, V\_k)} and \texttt{(V$_{1}$', V$_{2}$', \dots{}, V\_k')} are two colourings, then
\texttt{(V$_{1}$, V$_{2}$, \dots{}, V\_k) = (V$_{1}$', V$_{2}$', \dots{}, V\_k')} if and only if \texttt{V$_{i}$ = V$_{i}$'} for
\texttt{1 $\le$ i $\le$ k}.  For example, a triangle has the six distinct 3-colourings shown in
figure 8.7.  Note that even though there is exactly one vertex of each colour in
each colouring, we still regard these six colourings as distinct.}

The book then records the two anchoring cases: \emph{If \texttt{G} is empty, then each
vertex can be independently assigned any one of the \texttt{k} available colours.
Therefore \texttt{$\pi$\_k(G) = k\textasciicircum{}$\nu$}. On the other hand, if
\texttt{G} is complete, then there are \texttt{k} choices of colour for the first
vertex, \texttt{k-1} choices for the second, \texttt{k-2} for the third, and so on.
Thus, in this case, \texttt{$\pi$\_k(G) = k(k-1)\dots{}(k-$\nu$+1)}.}
\end{definition}

\begin{definition}[A ${u,v}$-component: $G(C \cup {u,v})$ for $C$ a component of $G - {u,v}$]
\lean{SimpleGraph.uvComponent}
\label{SimpleGraph.uvComponent}
\leanok
A \texttt{\{u,v\}}-component: \texttt{G[C $\cup$ \{u,v\}]} for \texttt{C} a component of
\texttt{G $-$ \{u,v\}}.

\emph{Let \texttt{S} be a vertex cut of a connected
graph \texttt{G}, and let the components of \texttt{G - S} have vertex sets
\texttt{V$_{1}$, V$_{2}$, \dots{}, V\_n}.
Then the subgraphs \texttt{G$_{i}$ = G[V$_{i}$ $\cup$ S]} are called the
\texttt{S}-components of \texttt{G} (see
figure 8.3). We say that colourings of \texttt{G$_{1}$, G$_{2}$, \dots{}, G\_n} agree on
\texttt{S} if, for
every \texttt{v $\in$ S}, vertex \texttt{v} is assigned the same colour in each of the
colourings.}

(The source prints \texttt{G$_{i}$ = G[V $\cup$ S]}; \texttt{V} there is an OCR slip for
\texttt{V$_{i}$}, as the
surrounding text and figure 8.3 make clear.)
\end{definition}

\begin{definition}[Type 1: every $(k-1)$-colouring assigns $u$ and $v$ the same colour]
\lean{SimpleGraph.Subgraph.IsType1}
\label{SimpleGraph.Subgraph.IsType1}
\leanok
Type 1: \textbf{every} \texttt{(k-1)}-colouring assigns \texttt{u} and \texttt{v} the
same colour.

\emph{We shall say that a \texttt{\{u, v\}}-component
\texttt{G$_{i}$} of \texttt{G} is of type 1 if every \texttt{(k-1)}-colouring of \texttt{G$_{i}$} assigns the same
colour to \texttt{u} and \texttt{v}} \dots{}
\end{definition}

\begin{definition}[Type 2: every $(k-1)$-colouring assigns $u$ and $v$ different colours]
\lean{SimpleGraph.Subgraph.IsType2}
\label{SimpleGraph.Subgraph.IsType2}
\leanok
Type 2: \textbf{every} \texttt{(k-1)}-colouring assigns \texttt{u} and \texttt{v}
different colours.

\dots{} \emph{and of type 2 if every
\texttt{(k-1)}-colouring of \texttt{G$_{i}$} assigns different colours to \texttt{u} and \texttt{v} (see figure
8.4).}
\end{definition}

\begin{definition}[$G$ contains a subdivision of $K_{4}$: four distinct branch vertices joined pairwise by six paths,\dots]
\lean{SimpleGraph.HasK4Subdivision}
\label{SimpleGraph.HasK4Subdivision}
\leanok
\texttt{G} contains a \textbf{subdivision of \texttt{K$_{4}$}}: four distinct branch vertices joined pairwise by six
paths, internally disjoint from each other and internally avoiding all four branch
vertices.

\emph{A subdivision of a graph \texttt{G} is a graph
that can be obtained from \texttt{G} by a sequence of edge subdivisions. A subdivision
of
\texttt{K$_{4}$} is shown in figure 8.5.}

The underlying operation is \S{}3.2: \emph{An edge \texttt{e} is said to be subdivided
when it is
deleted and replaced by a path of length two connecting its ends, the internal
vertex of this path being a new vertex.}

The book states Haj\'os' conjecture (1961) in the same paragraph: \emph{if \texttt{G} is
\texttt{k}-chromatic, then \texttt{G} contains a subdivision of \texttt{K\_k}} ---
adding that \emph{the condition is not sufficient; for example, a 4-cycle is a
subdivision of \texttt{K$_{3}$}, but is not 3-chromatic}, and that \texttt{k = 1, 2} are
obvious while \texttt{k = 3} holds \emph{because a 3-chromatic graph necessarily
contains an odd cycle, and every odd cycle is a subdivision of \texttt{K$_{3}$}}.
Theorem 8.5 (Dirac, 1952) settles \texttt{k = 4}.
\end{definition}

\begin{definition}[mycielskian]
\lean{SimpleGraph.mycielskian}
\label{SimpleGraph.mycielskian}
\leanok
\textbf{Mycielski's construction} on \texttt{V $\oplus$ V $\oplus$ Unit}: \texttt{inl} =
the old graph, \texttt{inr $\circ$ inl} = the
shadows \texttt{u$_{i}$} (joined to \texttt{v$_{i}$}'s neighbours), \texttt{inr $\circ$
inr} = the apex (joined to every shadow).

\emph{Suppose that we have
already constructed a triangle-free graph \texttt{G\_k} with chromatic number \texttt{k
$\ge$ 2}.
Let the vertices of \texttt{G\_k} be \texttt{v$_{1}$, v$_{2}$, \dots{}, v\_n}. Form a
new graph \texttt{G\_\{k+1\}} from
\texttt{G\_k} as follows: add \texttt{n+1} new vertices \texttt{u$_{1}$, u$_{2}$, \dots{}, u\_n, v}, and then, for
\texttt{1 $\le$ i $\le$ n}, join \texttt{u$_{i}$} to the neighbours of \texttt{v$_{i}$} and to \texttt{v}.  For example, if \texttt{G$_{2}$}
is \texttt{K$_{2}$} then \texttt{G$_{3}$} is the 5-cycle and \texttt{G$_{4}$} the
Gr\"otzsch graph (see figure 8.10).}
\end{definition}

\begin{definition}[The wheel with $n$ spokes: $C_n \lor K_{1}$]
\lean{SimpleGraph.wheel}
\label{SimpleGraph.wheel}
\leanok
\uses{SimpleGraph.join}
The \textbf{wheel} with \texttt{n} spokes: \texttt{C\_n $\lor$ K$_{1}$}.

(B\&M exercise 2.4.2, verbatim; used in exercise 8.4.5(b)). \emph{A wheel is a graph
obtained from a cycle by adding a new vertex and edges joining it to all the vertices of
the cycle; the new edges are called the spokes of the wheel.}
\end{definition}

\begin{definition}[$G$ is uniquely $k$-colourable: any two $k$-colourings induce the same partition]
\lean{SimpleGraph.UniquelyColorable}
\label{SimpleGraph.UniquelyColorable}
\leanok
\texttt{G} is \textbf{uniquely \texttt{k}-colourable}: any two \texttt{k}-colourings induce the same partition.

\emph{A graph \texttt{G} is uniquely
\texttt{k}-colourable if any two \texttt{k}-colourings of \texttt{G} induce the same partition of \texttt{V}.}
\end{definition}

\begin{definition}[$s$ is a maximal independent set of $G$ within $t$]
\lean{SimpleGraph.IsMaximalIndepSetIn}
\label{SimpleGraph.IsMaximalIndepSetIn}
\leanok
\texttt{s} is a \textbf{maximal independent set} of \texttt{G} within \texttt{t}.
\end{definition}

\begin{definition}[A canonical colouring, as a list of colour classes]
\lean{SimpleGraph.IsCanonicalColouring}
\label{SimpleGraph.IsCanonicalColouring}
\leanok
\uses{SimpleGraph.IsMaximalIndepSetIn}
A \textbf{canonical colouring}, as a list of colour classes.

\emph{A \texttt{k}-colouring \texttt{(V$_{1}$, V$_{2}$, \dots{}, V\_k)} of
\texttt{G} is said to be canonical if \texttt{V$_{1}$} is a maximal independent set of \texttt{G}, \texttt{V$_{2}$} is a
maximal independent set of \texttt{G - V$_{1}$}, \texttt{V$_{3}$} is a maximal
independent set of
\texttt{G - (V$_{1}$ $\cup$ V$_{2}$)}, and so on.  It is easy to see (exercise 8.6.3) that if \texttt{G} is
\texttt{k}-colourable, then there exists a canonical \texttt{k}-colouring of \texttt{G}.  By repeatedly
using the above method for finding maximal independent sets, one can determine all
the canonical colourings of \texttt{G}.  The least number of colours used in such a
colouring is then the chromatic number of \texttt{G}.}
\end{definition}

\begin{definition}[The Mycielski tower $G_{2} = K_{2}, G_{3}, G_{4}, \dots$]
\lean{SimpleGraph.mycielskiTower}
\label{SimpleGraph.mycielskiTower}
\leanok
\uses{SimpleGraph.mycielskian}
The Mycielski \textbf{tower} \texttt{G$_{2}$ = K$_{2}$, G$_{3}$, G$_{4}$, \dots{}}.
\end{definition}

\begin{theorem}[Every $k$-chromatic graph contains a $k$-critical subgraph]
\lean{SimpleGraph.exists_isKCritical_subgraph}
\label{SimpleGraph.exists_isKCritical_subgraph}
\leanok
\uses{SimpleGraph.IsKCritical}
Let $G$ be a finite simple graph whose chromatic number is $\chi(G) = k$. Then $G$
contains a subgraph $H$ that is $k$-critical, that is, $H$ satisfies $\chi(H) = k$ and
$\chi(H') < \chi(H)$ for every proper subgraph $H' \subset H$.
\end{theorem}

\begin{theorem}[Critical graphs are connected]
\lean{SimpleGraph.IsCritical.connected}
\label{SimpleGraph.IsCritical.connected}
\leanok
\uses{SimpleGraph.IsCritical}
Let $G$ be a graph on a finite, nonempty vertex set, and suppose $G$ is critical: every
proper subgraph $H \subset G$ satisfies $\chi(H) < \chi(G)$. Then $G$ is connected.
\end{theorem}

\begin{theorem}[Minimum degree of a $k$-critical graph]
\lean{SimpleGraph.IsKCritical.minDegree_ge}
\label{SimpleGraph.IsKCritical.minDegree_ge}
\leanok
\uses{SimpleGraph.IsKCritical}
Let $G$ be a finite $k$-critical graph, that is, a graph with chromatic number $\chi(G)
= k$ every proper subgraph of which has strictly smaller chromatic number. Then every
vertex of $G$ has degree at least $k - 1$; equivalently, the minimum degree satisfies
$\delta(G) \ge k - 1$.
\end{theorem}

\begin{theorem}[Lower bound on high-degree vertices in a k-chromatic graph]
\lean{SimpleGraph.card_filter_degree_ge_of_chromaticNumber}
\label{SimpleGraph.card_filter_degree_ge_of_chromaticNumber}
\leanok
Let $G$ be a finite simple graph whose chromatic number equals $k$. Then $G$ has at
least $k$ vertices of degree at least $k-1$; that is, the number of vertices $v$ with
$\deg_G(v) \ge k-1$ is at least $k$.
\end{theorem}

\begin{theorem}[Greedy bound on the chromatic number]
\lean{SimpleGraph.chromaticNumber_le_maxDegree_add_one}
\label{SimpleGraph.chromaticNumber_le_maxDegree_add_one}
\leanok
Let $G$ be a simple graph on a finite vertex set, and let $\Delta(G)$ denote its maximum
degree, that is, the largest number of neighbours of any vertex. Then the chromatic
number of $G$ satisfies $\chi(G) \le \Delta(G) + 1$.
\end{theorem}

\begin{theorem}[Vertex cuts of a critical graph are not cliques]
\lean{SimpleGraph.IsCritical.not_isClique_of_isVertexCut}
\label{SimpleGraph.IsCritical.not_isClique_of_isVertexCut}
\leanok
\uses{SimpleGraph.IsCritical, SimpleGraph.IsVertexCut}
Let $G$ be a finite connected graph that is critical, meaning that every proper subgraph
$H$ of $G$ satisfies $\chi(H) < \chi(G)$, where $\chi$ denotes the chromatic number.
Then no vertex cut of $G$ is a clique: if $S$ is a proper set of vertices whose deletion
leaves $G$ disconnected, then the vertices of $S$ are not pairwise adjacent.
\end{theorem}

\begin{theorem}[Critical graphs are connected and cut-vertex-free]
\lean{SimpleGraph.IsCritical.connected_and_no_cut_vertex}
\label{SimpleGraph.IsCritical.connected_and_no_cut_vertex}
\leanok
\uses{SimpleGraph.IsCritical, SimpleGraph.IsVertexCut}
Let $G$ be a critical graph on a finite, nonempty vertex set $V$. Then $G$ is connected,
and no single vertex is a vertex cut of $G$: for every $v \in V$, deleting $v$ leaves
the graph connected.
\end{theorem}

\begin{theorem}[Endpoints of a critical graph's 2-vertex cut are nonadjacent]
\lean{SimpleGraph.IsCritical.not_adj_of_isVertexCut_pair}
\label{SimpleGraph.IsCritical.not_adj_of_isVertexCut_pair}
\leanok
\uses{SimpleGraph.IsCritical, SimpleGraph.IsVertexCut}
Let $G$ be a finite simple graph that is critical, meaning that every proper subgraph of
$G$ has strictly smaller chromatic number than $G$. Suppose $u$ and $v$ are two distinct
vertices such that $\{u,v\}$ is a vertex cut of $G$, that is, deleting $u$ and $v$
leaves the remaining graph disconnected. Then $u$ and $v$ are not adjacent in $G$.
\end{theorem}

\begin{theorem}[Number of colors bounds the edge count]
\lean{SimpleGraph.chromaticNumber_ge_of_card_edges}
\label{SimpleGraph.chromaticNumber_ge_of_card_edges}
\leanok
Let $G$ be a finite simple graph on a nonempty vertex set, with $n$ vertices and $m$
edges. If $G$ admits a proper coloring using $k$ colors, then
\[
n^2 \le k\,\bigl(n^2 - 2m\bigr).
\]
\end{theorem}

\begin{theorem}[Five-colourability when odd cycles pairwise intersect]
\lean{SimpleGraph.chromaticNumber_le_five_of_odd_cycles_meet}
\label{SimpleGraph.chromaticNumber_le_five_of_odd_cycles_meet}
\leanok
Let $G$ be a finite simple graph. Suppose that every two odd cycles of $G$ share a
vertex; that is, for any two closed walks in $G$ that are cycles of odd length, there is
a vertex lying on both. Then the chromatic number of $G$ satisfies $\chi(G) \le 5$.
\end{theorem}

\begin{theorem}[Welsh–Powell bound on the chromatic number]
\lean{SimpleGraph.welsh_powell}
\label{SimpleGraph.welsh_powell}
\leanok
Let $G$ be a finite simple graph on $n$ vertices, and enumerate its vertices as $v_1,
v_2, \dots, v_n$ so that their degrees are arranged in non-increasing order, $d(v_1) \ge
d(v_2) \ge \dots \ge d(v_n)$. Then the chromatic number of $G$ satisfies
\[
\chi(G) \;\le\; \max_{1 \le i \le n} \min\bigl\{\, d(v_i) + 1,\; i \,\bigr\}.
\]
\end{theorem}

\begin{theorem}[A square-root bound on the chromatic number by edge count]
\lean{SimpleGraph.chromaticNumber_le_ceil_sqrt_two_mul_card_edges}
\label{SimpleGraph.chromaticNumber_le_ceil_sqrt_two_mul_card_edges}
\leanok
Let $G$ be a finite simple graph with $m \geq 1$ edges. Then the chromatic number of $G$
satisfies
\[
\chi(G) \le \left\lceil \sqrt{2m} \,\right\rceil.
\]
\end{theorem}

\begin{theorem}[Nordhaus–Gaddum upper bound on chromatic numbers]
\lean{SimpleGraph.nordhaus_gaddum}
\label{SimpleGraph.nordhaus_gaddum}
\leanok
Let $G$ be a simple graph on a finite vertex set $V$ with $n = \abs{V}$ vertices, and
let $G^{c}$ denote its complement, the graph on $V$ in which two distinct vertices are
adjacent precisely when they are non-adjacent in $G$. Writing $\chi$ for the chromatic
number, the sum of the chromatic numbers of a graph and its complement satisfies \[
\chi(G) + \chi(G^{c}) \le n + 1. \]
\end{theorem}

\begin{theorem}[The Szekeres–Wilf bound on chromatic number]
\lean{SimpleGraph.szekeres_wilf}
\label{SimpleGraph.szekeres_wilf}
\leanok
Let $G$ be a finite simple graph. For an induced subgraph $H$ of $G$, write $\delta(H)$
for its minimum degree (the least number of neighbours of any vertex of $H$). Then the
chromatic number of $G$ satisfies
\[
\chi(G) \le 1 + \max_{H} \delta(H),
\]
where the maximum is taken over all induced subgraphs $H$ of $G$.
\end{theorem}

\begin{theorem}[Gallai's two-per-class colouring theorem]
\lean{SimpleGraph.gallai_two_per_class}
\label{SimpleGraph.gallai_two_per_class}
\leanok
Let $G$ be a finite simple graph whose chromatic number equals $k$. Suppose $G$ admits a
proper colouring, using some number $n$ of colours, in which every one of the $n$ colour
classes contains at least two vertices. Then $G$ admits a proper colouring with exactly
$k$ colours in which, likewise, every colour class contains at least two vertices.
\end{theorem}

\begin{theorem}[Characterisation of 3-critical graphs as odd cycles]
\lean{SimpleGraph.isKCritical_three_iff_odd_cycle}
\label{SimpleGraph.isKCritical_three_iff_odd_cycle}
\leanok
\uses{SimpleGraph.IsKCritical}
Let $G$ be a finite simple graph on a nonempty vertex set. Then $G$ is $3$-critical if
and only if it is isomorphic to the cycle $C_n$ for some odd integer $n \ge 3$.
\end{theorem}

\begin{theorem}[The unique 1-critical graph is $K_1$]
\lean{SimpleGraph.isKCritical_one_iff}
\label{SimpleGraph.isKCritical_one_iff}
\leanok
\uses{SimpleGraph.IsKCritical}
Let $G$ be a finite graph on a nonempty vertex set. Then $G$ is $1$-critical if and only
if $G$ is isomorphic to $K_1$, the complete graph on a single vertex.
\end{theorem}

\begin{theorem}[The only 2-critical graph is $K_2$]
\lean{SimpleGraph.isKCritical_two_iff}
\label{SimpleGraph.isKCritical_two_iff}
\leanok
\uses{SimpleGraph.IsKCritical}
Let $G$ be a simple graph on a finite, nonempty vertex set. Then $G$ is $2$-critical if
and only if $G$ is isomorphic to the complete graph $K_2$ on two vertices.
\end{theorem}

\begin{theorem}[Vertex cuts of critical graphs are not uniquely colourable]
\lean{SimpleGraph.IsKCritical.not_uniquelyColorable_of_isVertexCut}
\label{SimpleGraph.IsKCritical.not_uniquelyColorable_of_isVertexCut}
\leanok
\uses{SimpleGraph.IsKCritical, SimpleGraph.IsVertexCut, SimpleGraph.UniquelyColorable}
Let $G$ be a finite simple graph on a nonempty vertex set, and let $k$ be a natural
number such that $G$ is $k$-critical (that is, $\chi(G) = k$ and $\chi(H) < k$ for every
proper subgraph $H$ of $G$). Let $S$ be a vertex cut of $G$, i.e.\ a proper subset of
the vertices whose deletion leaves $G$ disconnected. Then the subgraph of $G$ induced on
$S$ is not uniquely $(k-1)$-colourable: there exist two colourings of this induced
subgraph with $k-1$ colours that induce different partitions of $S$ into colour classes.
\end{theorem}

\begin{theorem}[Distinct vertices of a critical graph have incomparable neighbourhoods]
\lean{SimpleGraph.IsCritical.neighborSet_not_subset}
\label{SimpleGraph.IsCritical.neighborSet_not_subset}
\leanok
\uses{SimpleGraph.IsCritical}
Let $G$ be a critical graph on a finite, nonempty vertex set, so that $\chi(H) <
\chi(G)$ for every proper subgraph $H$ of $G$. Then for any two distinct vertices $u$
and $v$, the neighbourhood $N(u)$ is not contained in the neighbourhood $N(v)$.
\end{theorem}

\begin{theorem}[No $k$-critical graph has $k+1$ vertices]
\lean{SimpleGraph.no_isKCritical_card_eq_succ}
\label{SimpleGraph.no_isKCritical_card_eq_succ}
\leanok
\uses{SimpleGraph.IsKCritical}
Let $G$ be a finite graph on a nonempty vertex set $V$. If $G$ is $k$-critical — that
is, $G$ has chromatic number $\chi(G) = k$ and every proper subgraph $H$ of $G$
satisfies $\chi(H) < \chi(G)$ — then $G$ does not have exactly $k+1$ vertices; that is,
$\abs{V} \neq k+1$.
\end{theorem}

\begin{theorem}[Chromatic number of a join of graphs]
\lean{SimpleGraph.chromaticNumber_join}
\label{SimpleGraph.chromaticNumber_join}
\leanok
\uses{SimpleGraph.join}
Let $G$ and $H$ be simple graphs on finite, nonempty vertex sets, and let $G \vee H$
denote their join, the graph obtained from the disjoint union of $G$ and $H$ by adding
every edge between a vertex of $G$ and a vertex of $H$. Then the chromatic number of the
join is the sum of the chromatic numbers of the two factors: \[ \chi(G \vee H) = \chi(G)
+ \chi(H). \]
\end{theorem}

\begin{theorem}[Criticality of a join of graphs]
\lean{SimpleGraph.join_isCritical_iff}
\label{SimpleGraph.join_isCritical_iff}
\leanok
\uses{SimpleGraph.IsCritical, SimpleGraph.join}
Let $G$ and $H$ be simple graphs on finite, nonempty, disjoint vertex sets, and let $G
\vee H$ denote their join, the graph obtained from the disjoint union of $G$ and $H$ by
adding an edge between every vertex of $G$ and every vertex of $H$. Recall that a graph
is \emph{critical} when every proper subgraph has strictly smaller chromatic number than
the graph itself. Then $G \vee H$ is critical if and only if both $G$ and $H$ are
critical.
\end{theorem}

\begin{theorem}[The Hajós construction of $k$-critical graphs]
\lean{SimpleGraph.hajos_construction}
\label{SimpleGraph.hajos_construction}
\leanok
\uses{SimpleGraph.IsKCritical}
Let $G_1$ and $G_2$ be simple graphs on a common finite, non-empty vertex set, and
suppose that each becomes $k$-critical once its isolated vertices are discarded, i.e.
the subgraph induced on the non-isolated vertices of $G_i$ is $k$-critical for $i = 1,
2$. Suppose further that these two sets of non-isolated vertices have exactly one vertex
$v$ in common, and let $vv_1$ be an edge of $G_1$ and $vv_2$ an edge of $G_2$. Then the
graph
\[
\bigl((G_1 - vv_1) \cup (G_2 - vv_2)\bigr) + v_1v_2,
\]
obtained by deleting the edge $vv_1$ from $G_1$ and the edge $vv_2$ from $G_2$, forming
the union of the two resulting graphs, and adjoining the edge $v_1v_2$, is again
$k$-critical.
\end{theorem}

\begin{theorem}[Existence of 4-critical graphs on n vertices]
\lean{SimpleGraph.exists_isKCritical_four}
\label{SimpleGraph.exists_isKCritical_four}
\leanok
\uses{SimpleGraph.IsKCritical}
For every natural number $n$ with $n = 4$ or $n \ge 6$, there is a simple graph $G$ on
$n$ vertices that is $4$-critical: its chromatic number is $\chi(G) = 4$, and every
proper subgraph $H \subsetneq G$ satisfies $\chi(H) < \chi(G)$.
\end{theorem}

\begin{theorem}[Kainen's theorem on chromatic number and edge cuts]
\lean{SimpleGraph.kainen}
\label{SimpleGraph.kainen}
\leanok
Let $G$ be a finite graph, let $X$ be a set of its vertices with complement $V \setminus
X$, and let $n$ be a natural number. Suppose that the subgraphs induced on $X$ and on $V
\setminus X$ are each $n$-colourable, and that the edge cut between the two parts
contains fewer than $n$ edges — that is, the number of edges having one endpoint in $X$
and the other in $V \setminus X$ is at most $n-1$. Then $G$ is itself $n$-colourable.
\end{theorem}

\begin{theorem}[Edge connectivity of a critical graph]
\lean{SimpleGraph.IsKCritical.edgeConnectivity_ge}
\label{SimpleGraph.IsKCritical.edgeConnectivity_ge}
\leanok
\uses{SimpleGraph.IsKCritical, SimpleGraph.edgeConnectivity}
Let $G$ be a finite graph on a nonempty vertex set, and let $k$ be a natural number. If
$G$ is $k$-critical — that is, $G$ has chromatic number $k$ and every proper subgraph of
$G$ has strictly smaller chromatic number — then the edge connectivity of $G$ satisfies
$\kappa'(G) \ge k - 1$; equivalently, $G$ is $(k-1)$-edge-connected.
\end{theorem}

\begin{theorem}[Brooks' theorem on chromatic number and maximum degree]
\lean{SimpleGraph.brooks_chromaticNumber_le_maxDegree}
\label{SimpleGraph.brooks_chromaticNumber_le_maxDegree}
\leanok
Let $G$ be a connected simple graph on a finite, nonempty vertex set that is neither a
complete graph nor isomorphic to a cycle of odd length. Then the chromatic number of $G$
is at most its maximum degree, $\chi(G) \le \Delta(G)$.
\end{theorem}

\begin{theorem}[Equivalence of Brooks' theorem and an edge bound for critical graphs]
\lean{SimpleGraph.brooks_iff_edge_bound}
\label{SimpleGraph.brooks_iff_edge_bound}
\leanok
\uses{SimpleGraph.IsKCritical}
The following two statements are equivalent. (i) (Brooks' theorem) Every finite nonempty
connected graph $G$ that is neither a complete graph nor an odd cycle satisfies $\chi(G)
\le \Delta(G)$, where $\chi(G)$ denotes the chromatic number and $\Delta(G)$ the maximum
degree. (ii) For every integer $k \ge 4$, every finite nonempty $k$-critical graph $G$
that is not complete satisfies $\nu(k-1) + 1 \le 2\varepsilon$, where $\nu$ is the
number of vertices and $\varepsilon$ the number of edges of $G$.
\end{theorem}

\begin{theorem}[Four-colouring the line graph of a subcubic graph]
\lean{SimpleGraph.chromaticIndex_le_four_of_maxDegree_three}
\label{SimpleGraph.chromaticIndex_le_four_of_maxDegree_three}
\leanok
Let $G$ be a finite simple graph whose maximum degree $\Delta(G)$ equals $3$. Then the
line graph of $G$ — the graph whose vertices are the edges of $G$, two being adjacent
when they share an endpoint — has chromatic number at most $4$.
\end{theorem}

\begin{theorem}[Dirac's structure theorem for $k$-critical graphs with a $2$-vertex cut]
\lean{SimpleGraph.kCritical_two_vertex_cut_structure}
\label{SimpleGraph.kCritical_two_vertex_cut_structure}
\leanok
\uses{SimpleGraph.IsKCritical, SimpleGraph.IsVertexCut, SimpleGraph.Subgraph.IsType1, SimpleGraph.Subgraph.IsType2, SimpleGraph.contractEdge, SimpleGraph.uvComponent}
(Dirac, 1953.) Let $G$ be a finite, nonempty $k$-critical graph, and let $u \neq v$ be
two vertices for which $\{u,v\}$ is a $2$-vertex cut of $G$. Then $G - \{u,v\}$ has two
distinct components $C_1$ and $C_2$ whose $\{u,v\}$-components $G_1 = G[C_1 \cup
\{u,v\}]$ and $G_2 = G[C_2 \cup \{u,v\}]$ together make up all of $G$, that is $G = G_1
\cup G_2$; here $G_1$ is of type $1$ (every $(k-1)$-colouring of $G_1$ assigns $u$ and
$v$ the same colour) and $G_2$ is of type $2$ (every $(k-1)$-colouring of $G_2$ assigns
$u$ and $v$ different colours). Moreover, the graph $G_1 + uv$ obtained by adding the
edge $uv$ to $G_1$ and the graph $G_2 \cdot uv$ obtained by identifying $u$ and $v$ in
$G_2$ are both $k$-critical.
\end{theorem}

\begin{theorem}[Degree sum at a two-vertex cut in a k-critical graph]
\lean{SimpleGraph.kCritical_two_vertex_cut_degree_sum}
\label{SimpleGraph.kCritical_two_vertex_cut_degree_sum}
\leanok
\uses{SimpleGraph.IsKCritical, SimpleGraph.IsVertexCut}
Let $G$ be a finite, nonempty $k$-critical graph, and let $u,v$ be two distinct vertices
such that $\{u,v\}$ is a vertex cut of $G$, meaning that deleting $u$ and $v$ leaves the
remaining graph disconnected. Then the degrees of $u$ and $v$ satisfy \[ 3k \le d(u) +
d(v) + 5, \] equivalently $d(u) + d(v) \ge 3k - 5$.
\end{theorem}

\begin{theorem}[Dirac's theorem: 4-chromatic graphs contain a $K_4$-subdivision]
\lean{SimpleGraph.fourChromatic_hasK4Subdivision}
\label{SimpleGraph.fourChromatic_hasK4Subdivision}
\leanok
\uses{SimpleGraph.HasK4Subdivision}
Let $G$ be a finite simple graph whose chromatic number equals $4$. Then $G$ contains a
subdivision of $K_4$: there are four distinct branch vertices joined pairwise by six
paths that are internally disjoint from one another and internally avoid all four branch
vertices.
\end{theorem}

\begin{theorem}[Few low-degree vertices force a $K_4$-subdivision]
\lean{SimpleGraph.hasK4Subdivision_of_few_low_degree}
\label{SimpleGraph.hasK4Subdivision_of_few_low_degree}
\leanok
\uses{SimpleGraph.HasK4Subdivision}
Let $G$ be a finite simple graph on at least four vertices in which at most one vertex
has degree less than $3$. Then $G$ contains a subdivision of $K_4$.
\end{theorem}

\begin{theorem}[Edge count forcing a $K_4$-subdivision]
\lean{SimpleGraph.hasK4Subdivision_of_card_edges}
\label{SimpleGraph.hasK4Subdivision_of_card_edges}
\leanok
\uses{SimpleGraph.HasK4Subdivision}
Let $G$ be a finite simple graph with $n \ge 4$ vertices and at least $2n - 2$ edges.
Then $G$ contains a subdivision of $K_4$: four distinct branch vertices, joined in pairs
by six paths that are internally disjoint from one another and whose interiors avoid all
four branch vertices.
\end{theorem}

\begin{theorem}[Deletion–contraction recurrence for the chromatic polynomial]
\lean{SimpleGraph.numColorings_add_contract}
\label{SimpleGraph.numColorings_add_contract}
\leanok
\uses{SimpleGraph.contractEdge, SimpleGraph.numColorings}
Let $G$ be a finite simple graph, let $u$ and $v$ be adjacent vertices of $G$, and let
$k$ be a natural number. Writing $\pi_k(H)$ for the number of proper $k$-colourings of a
graph $H$, let $G - uv$ be the graph obtained from $G$ by deleting the edge $uv$, and
let $G \cdot uv$ be the graph obtained by contracting that edge, i.e. by identifying $u$
and $v$ into a single vertex. Then
\[
\pi_k(G) + \pi_k(G \cdot uv) = \pi_k(G - uv).
\]
\end{theorem}

\begin{theorem}[Existence of the chromatic polynomial]
\lean{SimpleGraph.exists_chromaticPolynomial}
\label{SimpleGraph.exists_chromaticPolynomial}
\leanok
\uses{SimpleGraph.numColorings}
Let $G$ be a finite simple graph on a nonempty vertex set, and write $n$ for the number
of its vertices and $\pi_k(G)$ for the number of distinct proper $k$-colourings of $G$.
Then there is a polynomial $p \in \bbz[x]$ with integer coefficients such that $p(k) =
\pi_k(G)$ for every $k \in \bbn$; moreover $p$ has degree $n$, leading coefficient $1$,
and constant term $0$, and its coefficients alternate in sign, in the sense that
$(-1)^{n-i}\,a_i \ge 0$ for every $i$, where $a_i$ denotes the coefficient of $x^i$ in
$p$.
\end{theorem}

\begin{theorem}[Second-highest coefficient of the chromatic polynomial]
\lean{SimpleGraph.chromaticPolynomial_coeff_card_sub_one}
\label{SimpleGraph.chromaticPolynomial_coeff_card_sub_one}
\leanok
\uses{SimpleGraph.numColorings}
Let $G$ be a finite simple graph on a nonempty vertex set $V$, with $\nu = |V|$ vertices
and $\varepsilon$ edges, and let $\pi_k(G)$ denote the number of proper $k$-colourings
of $G$. If $p \in \bbz[x]$ is a polynomial with integer coefficients satisfying $p(k) =
\pi_k(G)$ for every natural number $k$, then the coefficient of $x^{\nu - 1}$ in $p$
equals $-\varepsilon$.
\end{theorem}

\begin{theorem}[Chromatic polynomial of a tree]
\lean{SimpleGraph.numColorings_of_isTree}
\label{SimpleGraph.numColorings_of_isTree}
\leanok
\uses{SimpleGraph.numColorings}
Let $G$ be a tree on a finite, nonempty vertex set with $\nu$ vertices, and let $k$ be a
nonnegative integer. Then the number of distinct proper $k$-colourings of $G$ is
\[\pi_k(G) = k(k-1)^{\nu-1}.\]
\end{theorem}

\begin{theorem}[Chromatic polynomial bound for connected graphs]
\lean{SimpleGraph.numColorings_le_of_connected}
\label{SimpleGraph.numColorings_le_of_connected}
\leanok
\uses{SimpleGraph.numColorings}
Let $G$ be a connected simple graph on a finite, nonempty vertex set of size $\nu$, and
let $k$ be a natural number. Writing $\pi_k(G)$ for the number of distinct proper
$k$-colourings of $G$, one has
\[ \pi_k(G) \le k(k-1)^{\nu-1}, \]
and moreover, if $k \ge 3$, then equality $\pi_k(G) = k(k-1)^{\nu-1}$ holds if and only
if $G$ is a tree.
\end{theorem}

\begin{theorem}[Chromatic polynomial of the cycle graph]
\lean{SimpleGraph.numColorings_cycleGraph}
\label{SimpleGraph.numColorings_cycleGraph}
\leanok
\uses{SimpleGraph.numColorings}
Let $n \ge 3$, let $C_n$ denote the cycle graph on $n$ vertices, and let $k$ be a
nonnegative integer. Then the number of distinct proper $k$-colourings of $C_n$ is
\[
\pi_k(C_n) = (k-1)^n + (-1)^n\,(k-1).
\]
\end{theorem}

\begin{theorem}[Chromatic polynomial of the join with a single vertex]
\lean{SimpleGraph.numColorings_join_singleton}
\label{SimpleGraph.numColorings_join_singleton}
\leanok
\uses{SimpleGraph.join, SimpleGraph.numColorings}
Let $G$ be a finite graph, and let $\pi_k(H)$ denote the number of distinct proper
$k$-colourings of a graph $H$. Write $G \vee K_1$ for the join of $G$ with a single
vertex — the graph obtained from $G$ by adjoining one new vertex and joining it to every
vertex of $G$. Then for every integer $k \ge 1$,
\[
  \pi_k(G \vee K_1) = k\,\pi_{k-1}(G).
\]
\end{theorem}

\begin{theorem}[Chromatic polynomial of the wheel]
\lean{SimpleGraph.numColorings_wheel}
\label{SimpleGraph.numColorings_wheel}
\leanok
\uses{SimpleGraph.numColorings, SimpleGraph.wheel}
Let $n \ge 3$ and $k \ge 1$ be integers, and let $W_n = C_n \vee K_1$ be the wheel with
$n$ spokes, obtained from the cycle $C_n$ on $n$ vertices by adjoining a single hub
vertex joined to every vertex of the cycle. Writing $\pi_k(W_n)$ for the number of
distinct proper $k$-colourings of $W_n$, one has
\[ \pi_k(W_n) = k(k-2)^n + (-1)^n\, k(k-2). \]
\end{theorem}

\begin{theorem}[Chromatic count factors over connected components]
\lean{SimpleGraph.numColorings_eq_prod_components}
\label{SimpleGraph.numColorings_eq_prod_components}
\leanok
\uses{SimpleGraph.numColorings}
Let $G$ be a finite simple graph and $k$ a natural number, and write $\pi_k(G)$ for the
number of proper $k$-colourings of $G$. Then
\[
  \pi_k(G) = \prod_{C} \pi_k\bigl(G[C]\bigr),
\]
where the product ranges over the connected components $C$ of $G$ and $G[C]$ denotes the
subgraph of $G$ induced on the vertex set of $C$.
\end{theorem}

\begin{theorem}[A colouring-count identity for a union of subgraphs meeting in a clique]
\lean{SimpleGraph.numColorings_union_mul_inter}
\label{SimpleGraph.numColorings_union_mul_inter}
\leanok
\uses{SimpleGraph.numColorings}
Let $G$ be a simple graph on a finite vertex set, and let $G'$ and $H'$ be two subgraphs
of $G$; write $\pi_k(\cdot)$ for the number of distinct proper colourings of a graph
with $k$ available colours, and regard each of $G'$, $H'$, their join $G' \cup H'$, and
their meet $G' \cap H'$ as a graph on its own vertex set. Suppose the intersection $G'
\cap H'$ is complete, that is, every two of its vertices are adjacent in it. Then for
every $k \in \bbn$, \[ \pi_k(G' \cup H)\,\pi_k(G' \cap H') = \pi_k(G')\,\pi_k(H'), \]
where $G' \cup H'$ is written for the join.
\end{theorem}

\begin{theorem}[No root of the chromatic polynomial exceeds the vertex count]
\lean{SimpleGraph.chromaticPolynomial_no_root_gt_card}
\label{SimpleGraph.chromaticPolynomial_no_root_gt_card}
\leanok
\uses{SimpleGraph.numColorings}
Let $G$ be a nonempty finite simple graph on $\nu$ vertices, and for a natural number
$k$ write $\pi_k(G)$ for the number of proper $k$-colourings of $G$. Suppose $p$ is a
real polynomial that agrees with this count at every natural number, in the sense that
$p(k) = \pi_k(G)$ for all $k \in \bbn$. Then every real root of $p$ is at most $\nu$;
that is, if $p(x) = 0$ with $x \in \bbr$, then $x \le \nu$.
\end{theorem}

\begin{theorem}[Triangle-free graphs of every chromatic number]
\lean{SimpleGraph.exists_triangleFree_chromaticNumber_eq}
\label{SimpleGraph.exists_triangleFree_chromaticNumber_eq}
\leanok
For every positive integer $k$, there exists a finite simple graph $G$ that contains no
triangle and whose chromatic number is exactly $k$.
\end{theorem}

\begin{theorem}[$k$-criticality of the graphs in the Mycielski tower]
\lean{SimpleGraph.mycielskiTower_isKCritical}
\label{SimpleGraph.mycielskiTower_isKCritical}
\leanok
\uses{SimpleGraph.IsKCritical, SimpleGraph.mycielskiTower}
For every integer $k \ge 2$, the graph $G_k$ in the Mycielski tower is $k$-critical: its
chromatic number is $\chi(G_k) = k$, and it is critical, meaning that every proper
subgraph $H \subsetneq G_k$ satisfies $\chi(H) < \chi(G_k)$.
\end{theorem}

\begin{theorem}[Raising the chromatic number at girth at least six]
\lean{SimpleGraph.descartes_construction}
\label{SimpleGraph.descartes_construction}
\leanok
Let $k \ge 2$ be an integer, and let $G$ be a finite graph on a nonempty vertex set
whose chromatic number equals $k$ and whose girth is at least $6$. Then there exists a
finite graph $H$ whose chromatic number is at least $k+1$ and whose girth is at least
$6$.
\end{theorem}

\begin{theorem}[Graphs of high chromatic number and large girth]
\lean{SimpleGraph.exists_chromaticNumber_ge_girth_six}
\label{SimpleGraph.exists_chromaticNumber_ge_girth_six}
\leanok
For every integer $k \ge 2$, there exists a finite simple graph $H$ whose chromatic
number is at least $k$ and whose girth is at least $6$; that is, $H$ cannot be properly
colored with fewer than $k$ colors, yet every cycle in $H$ has length at least $6$.
\end{theorem}

\begin{theorem}[Existence of a canonical colouring for colourable graphs]
\lean{SimpleGraph.exists_isCanonicalColouring}
\label{SimpleGraph.exists_isCanonicalColouring}
\leanok
\uses{SimpleGraph.IsCanonicalColouring}
Let $G$ be a finite simple graph and let $k$ be a natural number. If $G$ is
$k$-colourable, then $G$ admits a canonical colouring consisting of at most $k$ colour
classes; that is, there is a list $(V_1, V_2, \dots, V_m)$ of pairwise disjoint sets of
vertices with $m \le k$, whose union is the whole vertex set, in which each $V_i$ is a
maximal independent set of $G$ among the vertices left uncoloured by $V_1, \dots,
V_{i-1}$.
\end{theorem}
